%
% einleitung.tex -- Beispiel-File für die Einleitung
%
% (c) 2020 Prof Dr Andreas Müller, Hochschule Rapperswil
%
% !TEX root = ../../buch.tex
% !TEX encoding = UTF-8
%
\section{Einführung\label{moebius:section:teil0}}
\kopfrechts{Möbius-Transformationen der komplexen Ebene}

Die komplexe Ebene stellt ein zentrales Objekt der komplexen Analysis dar. 
\index{komplexe Ebene}%
Punkte dieser Ebene können durch komplexe Zahlen beschrieben werden, wodurch algebraische und geometrische Strukturen miteinander verknüpft werden. 
Insbesondere ermöglichen Funktionen komplexer Variablen, Abbildungen der komplexen Ebene zu untersuchen und deren geometrische Eigenschaften zu analysieren.
Ein besonders wichtiges Thema solcher Abbildungen sind die sogenannten Möbius-Transformationen. 
\index{Mobius-Transformation@Möbius-Transformation}%
Diese Transformationen besitzen eine spezielle algebraische Struktur und zeigen bemerkenswerte geometrische Eigenschaften. 
Eine zentrale geometrische Eigenschaft von Möbius-Transformationen ist, dass sie Kreise und Geraden der komplexen Ebene wieder auf Kreise oder Geraden abbilden.
In der erweiterten komplexen Ebene werden Geraden dabei als spezielle Kreise betrachtet, die durch den Punkt im Unendlichen verlaufen. 
\index{Punkt im Unendlichen}%
Daher kann man sagen, dass Möbius-Transformationen sogenannte verallgemeinerte Kreise auf verallgemeinerte Kreise abbilden.



